% ==============================================================
\section{関連研究}\label{sec:related}
% ==============================================================

本研究は複数の研究領域にまたがる．
各領域をレビューし，本研究の貢献が埋めるギャップを明確にする．

\subsection{頻出パターンマイニングと密集区間マイニング}

Aprioriアルゴリズム~\cite{agrawal1994apriori}および
FP-Growth~\cite{han2000fpgrowth}は，頻出アイテムセットマイニングの基盤を確立した．
密集区間マイニングは，パターンのスライディングウィンドウサポートが閾値を超える
連続的な時間区間を特定することで，このフレームワークを拡張する．
密集区間は高活性期間を効果的に捉えるが，持続的な\emph{低活性}期間を
対称的に特徴づける機構を備えていない．
本研究の反密集区間の概念がこのギャップを埋める．

\subsection{希少・低頻度パターンマイニング}

低頻度または希少なパターンのマイニングは注目を集めている．
RP-Tree~\cite{szathmary2011rptree}はFP-Growthを拡張して
希少アイテムセットをマイニングする．
HaglinとManning~\cite{haglin2007minimal}は極小低頻度アイテムセットを導入した．
ERIM~\cite{akdas2024erim}は四つの希少アイテムセットマイニングアルゴリズムを
アンサンブルとして統合する．
データストリームにおける希少部分周期パターンに関する最近の研究~\cite{kiran2025r3pstream}は，
ストリーミング環境での時間的希少性に取り組んでいる．
負の相関ルール~\cite{wu2024negative}はアイテムの共欠如を捉えるが，
静的かつ非時間的な枠組みにとどまる．

これらのアプローチは\emph{大域的に}低頻度なパターンに焦点を当てているが，
低頻度の\emph{時間区間}を特徴づけるものではない．
パターンは大域的には頻出でありながら，持続的な不在期間を示すことがあり---
本研究の反密集区間はまさにこの現象を直接捉える．

\subsection{出現パターンとコントラストパターン}

DongとLi~\cite{dong1999emerging}は，二つのデータセット間のサポート比が
閾値を超えるアイテムセットとして出現パターン（EP）を導入した．
コントラストデータマイニングフレームワーク~\cite{dong2012contrast}は
これを多クラス比較に一般化した．
Garc\'ia-Vicoら~\cite{garciavico2018ep}はEPマイニング手法の包括的な分類体系を提供する．
Loyola-Gonz\'alezら~\cite{loyola2020contrast}は統計的検証を伴う
コントラストパターンマイニングをサーベイしている．
Liら~\cite{li2025hucspm}の最近の研究は，コントラストパターンを
高ユーティリティ系列パターンに拡張している．

出現パターンはデータセット間の\emph{サポート値}を比較する．
本研究のコントラスト密集パターンは時間レジーム間の\emph{密集区間構造}を
比較する---これはサポートが変化したかだけでなく，
\emph{サポートの時間的パターンがどのように変化したか}を捉える，
本質的に異なるより豊かな比較である．

\subsection{変化点検出}

変化点検出は，時系列の統計的性質が変化する時点を特定する．
CUSUM~\cite{page1954cusum}は平均シフトを逐次的に検出する．
PELT~\cite{killick2012pelt}は$O(n)$時間での最適な複数変化点検出を提供する．
BaiとPerron~\cite{bai2003structural}は回帰モデルにおける複数の構造変化に対処する．
Truongら~\cite{truong2020ruptures}はオフライン変化点検出のための
包括的なPythonライブラリを提供する．
最近の応用として金融時系列分析~\cite{chen2025pelt}がある．

これらの手法は\emph{単変量または多変量時系列}に対して動作し，
一般的な分布の変化を検出する．
\emph{パターンレベル}の密集区間の特定の構造を考慮せず，
本研究のフレームワークが行うような構造変化の\emph{種類}
（出現・消失等）の分類も行わない．

\subsection{概念ドリフト検出}

概念ドリフト検出は，時間とともにデータ生成過程が変化することを特定する．
ADWIN~\cite{bifet2007adwin}は統計的境界を用いた適応的ウィンドウ法を使用する．
DDM~\cite{gama2004ddm}はエラー率の監視によりドリフトを検出する．
ADWIN-U~\cite{assis2024adwinu}はADWINを教師なし環境に拡張する．
Cerqueiraら~\cite{cerqueira2024benchmark}は完全教師なしドリフト検出器の
ベンチマークを行っている．

概念ドリフト検出は「何かが変化した」というシグナルを提供するが，
個々のパターンレベルで\emph{何が}変化したかの構造化された特徴づけは提供しない．
本研究のトポロジー変化分類体系（出現・消失・増幅・収縮）は，
まさにこの解釈可能な特徴づけを提供する．

\subsection{統計的に有意なパターンマイニング}

パターンマイニングにおける統計的有意性の保証は重要な研究課題である．
Pellegrinaら~\cite{pellegrina2019westfallyoung}はWestfall--Young置換検定を
パターンマイニングに適用し，多重検定補正を伴う有意パターン発見のフレームワークを提示した．
Pellegrina and Vandin~\cite{pellegrina2020significant}は
置換検定に基づく最も有意なパターンの効率的マイニング手法を提案している．

これらの研究は\emph{静的}なパターン有意性を扱うのに対し，
本研究の置換検定はパターンの\emph{時間的構造変化}の有意性を評価する点で異なる．
ただし，多重検定補正の技法（Benjamini--Hochberg法）は共通の基盤である．

\subsection{時間パターンマイニングと周期パターンマイニング}

時間パターンマイニングは広く研究されてきた．
3P-ECLAT~\cite{3peclat2024}は列指向データベースにおける部分周期パターンをマイニングする．
TaTIRP~\cite{tatirp2025}はAllenの時間関係を用いて対象時間区間関連パターンを発見する．

これらのアプローチはイベント間の周期性や時間関係に焦点を当てており，
スライディングウィンドウサポート閾値により定義される
密集区間・反密集区間に焦点を当てる本研究とは相補的である．

\subsection{本研究の位置づけ}

表~\ref{tab:related_comparison}に本研究と関連研究の位置づけを要約する．
反密集区間は「低頻度の\emph{時間区間}」を捉える唯一のアプローチであり，
コントラスト密集パターンは密集区間\emph{構造}のレジーム間比較を
初めて形式化するものである．

\begin{table}[t]
\centering
\caption{関連研究との比較}
\label{tab:related_comparison}
\begin{tabular}{lcccc}
\toprule
手法カテゴリ & 時間的 & 構造比較 & 統計検定 & 低頻度 \\
\midrule
密集区間マイニング & \checkmark & -- & -- & -- \\
希少パターン & -- & -- & -- & \checkmark \\
出現パターン & -- & \checkmark & -- & -- \\
変化点検出 & \checkmark & -- & \checkmark & -- \\
概念ドリフト & \checkmark & -- & \checkmark & -- \\
有意パターン & -- & -- & \checkmark & -- \\
\textbf{反密集区間（本研究）} & \checkmark & -- & -- & \checkmark \\
\textbf{コントラスト密集（本研究）} & \checkmark & \checkmark & \checkmark & -- \\
\bottomrule
\end{tabular}
\end{table}
